\documentclass[11pt,a4paper]{article}
\usepackage[utf8]{inputenc}
\usepackage[T1]{fontenc}
\usepackage[portuguese]{babel}
\usepackage[a4paper,margin=2.5cm]{geometry}
\usepackage{booktabs}
\usepackage{amsmath}
\usepackage{amssymb}
\usepackage{verbatim}

\title{Segurança de Dados\\\large Assignment 5: Assinaturas digitais}
\author{Hugo Soares}
\date{18 de novembro de 2026}

\begin{document}
\maketitle

\section*{Questão 1}

\subsection*{Q1.a --- HMAC-SHA256 versus RSA-PSS}

\textbf{Ideia.} O HMAC é simétrico \cite{rfc2104}: as duas partes partilham
uma chave $K$ e calculam
\[
T = \mathrm{HMAC\text{-}SHA256}(K, M).
\]
A assinatura RSA-PSS é
assimétrica: o signatário usa a chave privada e o verificador usa a chave
pública,
$s = \mathrm{RSA}_{d}(\mathrm{EMSA\text{-}PSS}(H(M)))$ \cite{rfc8017}.

O programa assina e verifica a mesma mensagem de 37 bytes (uma transferência de
500 EUR para a conta PT50-0001) e mede cada operação
individualmente com \texttt{time.perf\_counter\_ns()}, depois de um
aquecimento. Foram feitas 1000 iterações para o HMAC e 500 para o RSA, que é
mais lento.

\begin{verbatim}
tag = hmac.HMAC(key, hashes.SHA256()); tag.update(M); tag.finalize()
sig = private_key.sign(M, padding.PSS(mgf=padding.MGF1(hashes.SHA256()),
                                      salt_length=32), hashes.SHA256())
public_key.verify(sig, M, pss, hashes.SHA256())
\end{verbatim}

\begin{center}
\begin{tabular}{lcc}
\toprule
 & HMAC-SHA256 & RSA-PSS-2048/SHA-256 \\
\midrule
Chave usada & 32 B (partilhada) & 1218 B privada / 294 B pública \\
Tag/assinatura & 32 B & 256 B \\
Geração & 1,59 $\mu$s & 596,6 $\mu$s \\
Verificação & 1,54 $\mu$s & 20,2 $\mu$s \\
\bottomrule
\end{tabular}
\end{center}

\textbf{Discussão.} O HMAC é cerca de 375 vezes mais rápido a gerar e 13 vezes
mais rápido a verificar; a assinatura é 8 vezes maior. A assinatura RSA é
dominada pela exponenciação modular com o expoente privado $d$; a verificação
usa $e=65537$, por isso é muito mais barata. Em compensação, o HMAC exige uma
chave secreta partilhada: com $n$ participantes são necessárias
$n(n-1)/2$ chaves e cada par tem de as distribuir por um canal seguro. A chave
pública RSA pode ser publicada e certificada, e uma única chave privada serve
para assinar para todos. O custo de gerar o par de chaves RSA (não medido
acima) é igualmente muito superior ao de gerar 32 bytes aleatórios.

\subsection*{Q1.b --- Porque é que o HMAC não resolve a disputa?}

No HMAC, $K$ é conhecida por Alice e por Bob. Logo, Alice consegue calcular
exatamente a mesma tag que Bob para qualquer mensagem. Se Alice afirmar que
Bob lhe enviou $M$, Bob pode responder que foi Alice que gerou a tag, e um
juiz não tem forma de distinguir os dois: a tag apenas prova que
\emph{alguém que conhece $K$} criou a mensagem. O HMAC dá integridade e
autenticidade entre as partes que confiam uma na outra, mas não dá
\textbf{não-repúdio}, porque o verificador tem o mesmo poder de forjar que o
signatário.

Numa assinatura digital, o valor $s$ depende de $d$, que só Bob conhece. Alice
verifica com $e$, que é público, mas não consegue produzir $s$. É esta
assimetria de conhecimento que sustenta o não-repúdio: o requisito é que a
assinatura use \emph{apenas dados conhecidos pelo signatário}. Bob não pode
negar depois ter assinado, desde que a chave privada esteja protegida; se a
chave for comprometida, a garantia desaparece. Na prática isto exige também
uma PKI que ligue a chave pública à identidade de Bob, e revogação em caso de
compromisso.

\subsection*{Q1.c --- Cifrar com a chave pública de Alice não autentica a origem}

Cifrar com $PU_{Alice}$ garante confidencialidade para Alice: só ela tem
$PR_{Alice}$ para decifrar. Mas $PU_{Alice}$ é pública, portanto
\emph{qualquer pessoa} pode construir um ciphertext válido. O seguinte código
não usa Bob em lado nenhum:

\begin{verbatim}
forged = b"TRANSFER 99999 EUR TO ACCOUNT PT50-6666 -- Bob"
ct = alice_public.encrypt(forged, OAEP)
pt = alice_private.decrypt(ct, OAEP)   # Alice decifra sem erros
\end{verbatim}

A execução mostra que Alice obtém a mensagem forjada e não tem qualquer
indicação de que não foi Bob:

\begin{verbatim}
Alice decrypts successfully:
  TRANSFER 99999 EUR TO ACCOUNT PT50-6666 -- Bob
Mallory's signature does NOT verify with Bob's public key.
Bob's signature verifies with Bob's public key: the source is proven.
\end{verbatim}

O RSA-OAEP protege contra modificação do ciphertext, mas não diz quem o
criou. Para autenticar a origem é preciso uma assinatura: Bob assina com
$PR_{Bob}$ e Alice verifica com $PU_{Bob}$. Se também for necessário sigilo,
faz-se \emph{sign-then-encrypt} (assinar e depois cifrar para Alice). A
identidade escrita no plaintext não vale nada sem essa prova criptográfica.

\section*{Questão 2}

\subsection*{Q2.a --- Assinatura RSA de livro}

Com as chaves geradas pela biblioteca, obtêm-se $n$, $e$ e $d$ através de
\texttt{private\_numbers()} e \texttt{public\_numbers()}. Cada mensagem é
interpretada como um inteiro big-endian, e assina-se diretamente
$s = m^d \bmod n$, sem hash nem padding:

\begin{verbatim}
s = pow(m, d, n)
pow(s, e, n) == m          # verificação
\end{verbatim}

Para as duas mensagens do enunciado (295 bits cada), a assinatura e a
verificação funcionam, e alterar a mensagem faz a verificação falhar. O
esquema é determinístico: a mesma mensagem produz sempre a mesma assinatura.

\subsection*{Q2.b --- Forjar uma assinatura a partir de duas válidas}

Se $s_1 = m_1^d \bmod n$ e $s_2 = m_2^d \bmod n$, então
\[
s_1 s_2 \equiv (m_1 m_2)^d \pmod n .
\]
Basta portanto definir $m_3 = m_1 m_2 \bmod n$ e
$s_3 = s_1 s_2 \bmod n$; a assinatura $s_3$ é válida para $m_3$ sem que o
atacante conheça $d$:

\begin{verbatim}
m3 = (m1 * m2) % n
s3 = (s1 * s2) % n
pow(s3, e, n) == m3        # True, sem conhecer d
\end{verbatim}

Na execução, $m_1 m_2$ tem 581 bits (logo não há redução módulo $n$), e
$s_3$ verifica corretamente. A mensagem $m_3$ em bytes não é texto legível:
\[
\texttt{1b c6 15 49 4d 7c 54 c2 d6 db 23 5a 56 9b 8c f9 53 9e 79 91 4b 6b c1 bd \ldots}
\]
A forja é \textbf{existencial} (\emph{new message for an existing
signature}): produz uma \emph{mensagem nova} para assinaturas já existentes,
não uma assinatura fraudulenta para uma mensagem escolhida pelo atacante.
$m_3$ fica determinado pelo produto de $m_1$ e $m_2$; para um $m_3$ escolhido
a priori, o atacante teria de o fatorizar em $m_1 m_2$ e obter assinaturas
nesses fatores, pelo que esta construção não dá uma falsificação para uma
mensagem arbitrária. A falha mostra que a assinatura tem de resistir a forjas
existenciais, não apenas a falsificações de mensagens pré-escolhidas.

\subsection*{Q2.c --- Porque é que $E(PR_a, H(M))$ com PSS resiste?}

Com hash, o que é assinado é $H(M)$, e não $M$. O hash destrói a estrutura
multiplicativa: em geral
$H(m_1 m_2) \neq H(m_1)H(m_2) \bmod n$, e o produto de duas assinaturas deixa
de corresponder a uma mensagem com significado. Além disso, o PSS não assina
$H(M)$ directamente: constrói
$\mathrm{EMSA\text{-}PSS}(H(M))$ com um \emph{salt} aleatório, MGF1 e um
trailer fixo, e só depois aplica a exponenciação privada. Duas assinaturas da
mesma mensagem são diferentes, e a verificação descodifica a estrutura PSS e
rejeita qualquer resultado que não tenha o padding correcto. Para forjar seria
necessário inverter o RSA ou encontrar uma colisão/segunda pré-imagem do hash,
o que é computacionalmente inviável com SHA-256. O hash também resolve o
problema de tamanho (mensagens maiores que $n$) e liga a assinatura a toda a
mensagem.

\section*{Questão 3}

\subsection*{Q3.a --- Verificação DSA}

A verificação DSA \cite{fips186} calcula
\[
w = s^{-1} \bmod q,\quad u_1 = H(M)w \bmod q,\quad u_2 = rw \bmod q,
\]
\[
v = \left(g^{u_1} y^{u_2} \bmod p\right) \bmod q,
\]
e aceita se $v = r$. O programa \texttt{q3a\_dsa\_verify.py} implementa esta
verificação e confirma as duas assinaturas:

\begin{verbatim}
p: 1024 bits, q: 160 bits
Q3.a signature 1 valid: True
Q3.a signature 2 valid: True
\end{verbatim}

\subsection*{Q3.b --- O que revela um $r$ partilhado}

Como $r = (g^k \bmod p) \bmod q$, o mesmo $r$ nas duas assinaturas significa
que foi reutilizado o mesmo $k$. Isto revela que o signatário não gerou um
nonce novo por assinatura --- por exemplo, usou um $k$ fixo ou um gerador de
números aleatórios fraco. Da equação de assinatura
$s = k^{-1}(H(M) + xr) \bmod q$:
\[
s_1 = k^{-1}(H_1 + xr), \qquad s_2 = k^{-1}(H_2 + xr).
\]
Subtraindo,
\[
s_1 - s_2 = k^{-1}(H_1 - H_2) \pmod q
\quad\Longrightarrow\quad
k = (H_1 - H_2)(s_1 - s_2)^{-1} \bmod q .
\]
Conhecido $k$, obtém-se a chave privada:
\[
x = (s_1 k - H_1) r^{-1} \bmod q .
\]
Não é preciso resolver logaritmos discretos: a reutilização transforma o
problema num sistema linear em $\mathbb{Z}_q$. O mesmo acontece se $k$ for
previsível ou tiver pouca entropia.

\subsection*{Q3.c --- Recuperar $x$ e forjar uma assinatura}

O programa \texttt{q3c\_dsa\_attack.py} aplica as expressões acima com
$H(M) = \mathrm{int}(\mathrm{SHA\text{-}1}(M)) \bmod q$ e verifica o resultado
contra a chave pública:

\begin{verbatim}
Q3.b
  k  = 0x57981d4fe83ab4095a21184218b0971a379fa99a
  x  = 0x445240674519b2c71838654fdd1c1c9da2078c62
  y == g^x mod p: True

Q3.c
  message = 'TRANSFER 99999 EUR TO ACCOUNT PT50-6666'
  k3 = 0x5eed1234abcd0001
  r3 = 0x52de62fab535ceeadaf2bf0b8dc824c1ea8b460a
  s3 = 0x772474a8ce404dfd66b3a5a4a312dab61bf30875
  forged signature verifies: True
\end{verbatim}

A assinatura forjada, com um $k_3$ novo escolhido pelo atacante, passa a
verificação de Q3.a. A chave privada recuperada permite assinar qualquer
mensagem, pelo que a reutilização de $k$ compromete a chave por completo. Na
prática, cada assinatura deve usar um $k$ único e imprevisível, gerado por um
CSPRNG; a alternativa é a geração determinística do RFC 6979 \cite{rfc6979},
que deriva $k$ da chave e da mensagem. Note-se ainda que os parâmetros do
enunciado são DSA legado ($p$ de 1024 bits, $q$ de 160 bits, SHA-1); as
recomendações atuais pedem grupos maiores e SHA-2/SHA-3, ou ECDSA/EdDSA.

\section*{Conclusão}

O HMAC é a escolha eficiente para integridade e autenticidade entre partes que
partilham uma chave, mas não serve para não-repúdio, porque ambas conseguem
forjar a tag. A assinatura digital resolve a disputa porque assenta num segredo
que só o signatário possui. Cifrar com a chave pública de Alice dá
confidencialidade, não autenticação da origem: qualquer um pode cifrar. O RSA
de livro é maleável e permite forjas existenciais por multiplicação, o que o
hash e o padding PSS eliminam. Por fim, a reutilização do nonce em DSA revela
$k$ e $x$ por álgebra linear módulo $q$, o que permite assinar qualquer
mensagem; o nonce tem de ser único por assinatura.

\section*{Como executar}

\begin{verbatim}
python3 q1a_hmac_vs_rsa.py
python3 q1c_encrypt_to_alice.py
python3 q2_textbook_rsa.py
python3 q3a_dsa_verify.py
python3 q3c_dsa_attack.py
pdflatex report.tex && pdflatex report.tex
\end{verbatim}

\begin{thebibliography}{9}

\bibitem{rfc2104}
H. Krawczyk, M. Bellare, R. Canetti.
\newblock \emph{HMAC: Keyed-Hashing for Message Authentication}, RFC 2104, 1997.

\bibitem{rfc8017}
K. Moriarty, B. Kaliski, J. Jonsson, A. Rusch.
\newblock \emph{PKCS \#1: RSA Cryptography Specifications Version 2.2}, RFC 8017, 2016.

\bibitem{fips186}
NIST.
\newblock \emph{Digital Signature Standard (DSS)}, FIPS 186-5, 2023.

\bibitem{rfc6979}
T. Pornin.
\newblock \emph{Deterministic Usage of the Digital Signature Algorithm (DSA) and
Elliptic Curve Digital Signature Algorithm (ECDSA)}, RFC 6979, 2013.

\end{thebibliography}

\end{document}
